\subsection{Universal short products over arbitrary fields}

\begin{lemma}[Proved: unused long arrows split off two-spheres]
\label{A7:unused-long-wedges}
Let $X$ be a finite acyclic category with full support
$0<1<2<3<4$ and a singleton $X(0,4)$.
Delete every arrow in $X(0,3)$ and $X(1,4)$ which is not
a composite of two nonidentity arrows, obtaining a
subcategory $Y$ on the same objects. Put
$A=X(0,1)$, $D=X(3,4)$, and let $r,s$ be the respective
numbers of deleted $0\to3$ and $1\to4$ arrows. Then
\[
 BX\simeq BY\vee\bigvee^{\,r(|D|-1)+s(|A|-1)}S^2.
\]
This equivalence can be chosen under the specified
inclusion $BY\hookrightarrow BX$. In particular, vanishing
of all degree-two products on $BY$ implies the same
vanishing on $BX$, over any field.
\end{lemma}
\begin{proof}
First, every full-support subcategory retaining the
singleton extreme hom-set is simply connected. Indeed,
choose arrows $a_i:0\to i$, with $a_0=1_0$ and $a_4$
the unique extreme arrow. For $f:i\to j$, choose
$b_j:j\to4$, using $b_4=1_4$. In the universal groupoid,
\[
 b_j f a_i=a_4=b_j a_j.
\]
Cancellation gives $f=a_j a_i^{-1}$. Thus all arrows are
determined by a tree of the $a_i$, and the fundamental
group at $0$ is trivial. The assertion applies to $Y$:
each long hom-set retains a composite and hence remains
nonempty, and deletion of arrows not hit by composition
is closed under composition.

Adjoin one deleted arrow $e:0\to3$ to any such
subcategory. The only new nondegenerate simplices are
the edge $e$ and the triangles $(e,d)$ for $d\in D$.
No new simplex has dimension at least three, since
there is no nonidentity arrow into $0$ and only one
further object after $3$. Nor is a new triangle a face
of an old three-simplex: such a face would express
$e$ as a two-arrow composite, contrary to its definition.
Every composite $de$ is the unique extreme arrow.

Choose $d_0\in D$. The subcomplex obtained by adjoining
only $e$ and $(e,d_0)$ is an elementary expansion of
the old nerve realization: the new edge is a free
face of that triangle, and its other two edges already
belong to the old complex. It deformation retracts
onto the old complex while fixing it. Each of the
remaining $|D|-1$ triangles is a two-cell attached
along a loop in this subcomplex. That subcomplex is
simply connected by the first paragraph, so every
attaching map is null-homotopic. Attaching these
two-cells therefore gives a wedge of $|D|-1$ copies
of $S^2$, up to homotopy under the old complex.

Repeat for every unused $0\to3$ arrow. The dual
argument for an unused $1\to4$ arrow adjoins one
edge and $|A|$ triangles, giving $|A|-1$ two-spheres.
The two kinds of new arrows cannot be consecutive
factors, and neither changes $A$ or $D$. The asserted
formula follows. In a wedge with two-spheres every
positive product involving a sphere summand is zero,
while products of classes from $BY$ are their original
products. This proves the ring assertion over any field.
\end{proof}

\begin{theorem}[Proved: universal product transfer over any field]
\label{A7:universal-short-all-fields}
Use the three short maps and the universal category
$F=F(P,Q,R)$ of Theorem~\ref{A7:universal-short-model}.
For every field $k$, if all products
$H^2(BF;k)\otimes H^2(BF;k)\to H^4(BF;k)$ are zero,
then all such products on $B\mathcal C$ are zero.
\end{theorem}
\begin{proof}
Collapse $\mathcal C(0,4)$ to a singleton, obtaining
$\mathcal C_0$. Proposition~\ref{A7:extreme-arrow-collapse}
gives the actual product-preserving surjection from
its second cohomology to that of $\mathcal C$.
Delete its unused $03$ and $14$ arrows to obtain $Y$.
Lemma~\ref{A7:unused-long-wedges} shows that vanishing
of products transfers from $BY$ to $B\mathcal C_0$.

The same explicit functor $F\to Y$ used in
Theorem~\ref{A7:universal-short-model} is the identity
on the normalized chain groups in degrees three and
four and on the fourth boundary. Its map on fourth
cohomology is therefore an isomorphism over $k$.
For $u,v\in H^2(BY;k)$, naturality sends $uv$ to
the product of their two pullbacks on $BF$. The
hypothesis makes this zero, and the fourth-cohomology
isomorphism forces $uv=0$. Transfer successively
through the wedge decomposition and the degree-two
surjection proves the claim. No bound on second
Betti number and no characteristic-two argument
is required.
\end{proof}

\begin{corollary}[Proved: two components suffice for tree-stalk prefixes]
\label{A7:partial-long-tree}
In a partial initial assignment of the three short
maps $P,Q,R$, retain all vertices of the two known-edge
graphs from Proposition~\ref{A7:partial-long-graph}.
Every completion has zero rational degree-two products
if either of the following conditions holds:
\begin{enumerate}
\item the $03$ graph has at most two connected components,
and either $|U|\leq2$ or $|A|\leq2$, $|B|\leq3$;
\item the $14$ graph has at most two connected components,
and either $|W|\leq2$ or $|D|\leq2$, $|C|\leq3$.
\end{enumerate}
\end{corollary}
\begin{proof}
Adding the still unknown edges can only decrease the
number of connected components. In the first case,
every completed universal category has $|F(0,3)|\leq2$.
Its source stalks satisfy (S1) or (S2) of
Proposition~\ref{A7:tree-stalk-profile}, respectively.
Thus all its rational degree-two products are zero.
Theorem~\ref{A7:universal-short-all-fields} transfers
this to every category with those completed short
maps. The second case uses (T1) or (T2) at the sink.
All vertices, including unused short-map values,
are essential to the component count.
\end{proof}
